\section{논의}
\label{sec:discussion}

%======================================================================
\subsection{단계 위험을 좌우하는 요인}
\label{sec:discussion:drivers}

% P1: 구조보다 정식화와 맥락; 먼저 정식화
다음 단계 탈락 위험의 예측을 좌우한 것은 모델 구조보다 무엇을 추정하도록 학습하는가, 즉 정식화와 입력이 담는 맥락이었다.
먼저 상대위험 목적은 단계마다 공통적인 위험 수준을 배우지 못해 단계별 예측이 평평해지지만(\ref{sec:method:rq2}절, 그림~\ref{fig:contrast}(a)), 목적을 바꾼 이득은 같은 경기·같은 단계의 팀 사이에서도 남았다.
이는 다음 단계라는 짧은 구간에 맞춘 목적이 남은 경기 전체의 탈락 순서를 맞추는 목적과 다른 정보에 무게를 두기 때문으로 해석된다.
KRAFTON 적용판과 Cox 목적이 원래 과제에서는 앞선 결과(\ref{sec:results:rq2}절)도 학습 목적이 곧 추정 대상을 정함을 보이므로, 설계의 비교는 성능 순위보다 추정 대상이 의사결정의 시간 단위와 맞는지를 기준으로 읽어야 한다.

% P2: 입력은 특징 수보다 구역에 대한 처지; 구조의 이득이 작은 이유
입력 표현에서는 특징의 수보다 줄어드는 구역에 대한 팀의 처지를 담았는지가 중요하였다.
활동 점수는 플레이어가 무엇을 했는지를 요약할 뿐 구역과 주변 팀에 대해 어디에 있는지는 담지 않으며, 서식지 맥락 블록이 Sagong 설계와의 차이의 절반 이상을 설명하였다.
구역까지의 거리와 주변 팀 수를 이미 담은 KRAFTON 입력은 본 연구의 입력과 예측력이 같았지만 둘을 합치면 성능이 올랐는데, 무기별 피해와 장비 수준은 KRAFTON 입력에만, 이동 부담과 관계 요약은 본 연구의 입력에만 있다는 점에서 두 입력은 서로 다른 측면을 담는 것으로 보인다.
모델 구조의 이득이 작았던 것은 팀 사이의 관계가 주변 경쟁 밀도와 관계 요약 특징으로 이미 입력에 요약되어 있어, attention이 새로 찾아낼 정보가 적었기 때문일 수 있다.

% P3: 생태학적 해석
서식지 관점은 위험 요인의 효과가 단계에 따라 달라지는 양상을 일관되게 설명한다.
구역이 넓은 초반에는 주변 팀의 수가 위험을 크게 가르지만, 구역이 좁아진 후반에는 모든 팀이 가까이 모여 그 차이가 덜 중요해지고 남은 구역 안에서 얼마나 깊이 자리 잡았는지와 남은 팀원 수가 더 중요해지는 것으로 해석된다.
경계에 가까울수록 위험이 크고 그 차이가 구역이 작아질수록 커지는 것은, 서식지가 줄고 조각날수록 가장자리의 개체가 더 큰 위험에 놓이는 가장자리 효과와 같은 방향이다~\cite{Fahrig03}.
네 변수가 판별력의 대부분을 설명했다는 점에서, 이 관점은 비유를 넘어 위험을 요약하는 실용적인 틀이다.

%======================================================================
\subsection{활용과 시사점}
\label{sec:discussion:implications}

% P4: 경기 분석, 중계, 구역 설계
단계마다 보정된 위험은 경기 분석, 중계, 구역 설계에 같은 척도로 쓰일 수 있다.
단계별 평균 예측이 실제 탈락률과 맞아 단계 2의 30\%와 단계 6의 30\%가 같은 뜻을 가지므로, 코치와 분석가는 단계 경계마다 팀의 위험을 비교하며 결정을 복기하고, 중계는 다음 단계에 탈락할 가능성이 높은 팀을 미리 짚으며, 개발사는 구역 규칙 변경 전후의 단계별 탈락 양상을 비교할 수 있다.

% P5: 대규모 시공간 위험 예측에 대한 시사점
이 결과는 이동하는 다수의 에이전트를 다루는 대규모 시공간 위험 예측에도 세 가지 시사점을 준다.
첫째, 입력뿐 아니라 생존 판정과 에이전트 간 관계에도 같은 절단 시각을 적용하고, 절단 이후 기록을 지운 재구성으로 누수가 없음을 확인해야 한다.
둘째, 통합 AUC는 시점별 위험 수준을 맞히는 효과와 같은 상황의 개체를 가려내는 효과를 섞으므로, 단계 안 지표와 단계별 보정을 함께 보고해야 한다.
셋째, 설계를 비교할 때는 구조를 바꾸기 전에 추정 대상과 입력을 통제한 분해로 차이의 출처부터 확인해야 한다.

%======================================================================
\subsection{한계}
\label{sec:discussion:limitations}

% P6: 한계
본 연구의 결론은 다음 범위 안에서 읽어야 한다.
첫째, 입력은 모든 팀의 위치를 아는 관전자 관점에서 만들었으므로, 각 팀이 당시 알던 정보만 쓰는 의사결정 지원과는 구분된다.
둘째, 자료는 한 구역 규칙 아래 약 10주간의 스팀 서버 경기이고 솔로 코호트는 하나뿐이어서, 다른 규칙과 플랫폼, 솔로 경기로의 일반화는 추가 확인이 필요하다.
셋째, 두 기존 설계는 원래 자료와 조건을 재현한 것이 아니라 같은 자료에 맞춘 적용판이므로, 원 논문의 수치와 직접 비교할 수 없다.
넷째, RQ1의 계수는 관측 자료의 조건부 연관이며, 구역 밖 깊이와 이동 부족분처럼 상관된 요인의 효과는 분리하기 어렵다.
다섯째, 구조의 이득이 작다는 결론은 현재의 데이터 규모와 요약 특징을 쓰는 조건의 결과로, 더 많은 경기나 플레이어 수준의 궤적을 쓰면 달라질 수 있다.

%======================================================================
\section{결론}
\label{sec:conclusion}

% P7: 결론과 향후 과제
본 연구는 배틀로얄 경기의 대규모 텔레메트리에서 다음 단계 탈락 위험을 서식지 축소 관점의 단계 조건부 이산시간 hazard로 정식화하고, 누수 없는 팀--단계 표본과 통제된 분해로 이 위험의 예측을 무엇이 좌우하는지 분석하였다.
경쟁 밀도, 이동 부담, 남은 팀원 수의 효과는 코호트에 걸쳐 일관되었고 경쟁 밀도와 구역 안쪽 깊이의 효과는 단계에 따라 달라졌으며, 기존 설계와의 차이는 주로 정식화와 서식지 맥락에서 왔고 모델 구조의 기여는 작았다.
향후에는 각 팀이 관측할 수 있는 정보만 쓰는 의사결정 관점으로 모델을 확장하고, 수집 이후 적용된 구역 규칙 개편~\cite{PUBGBlueZone2026}을 자연 실험으로 삼아 위험 요인이 규칙이 바뀌어도 유지되는지 검증할 계획이다.

\section*{사사}
분석 코드 작성과 원고 초안 작성 및 교정에 생성형 AI(Anthropic Claude)를 사용하였으며, 저자가 모든 코드, 수치, 문장을 검토하고 수정하였다.
